\section{Pendahuluan}

\subsection{Latar Belakang}
Perangkat lunak modern jarang dibangun dari nol. Sebuah rilis biasanya dirakit dari kode
sendiri, pustaka pihak ketiga, dan artefak yang dibangun otomatis di pipeline CI/CD, lalu
diunggah ke \textit{registry} dan diunduh oleh pengguna. Setiap tahap dalam rantai ini dapat menjadi
titik masuk bagi perubahan yang tidak diinginkan, baik karena kesalahan maupun serangan. Risiko
ini dikenal sebagai risiko rantai pasok (\textit{supply chain}). Pengguna yang menerima artefak
perlu cara untuk memastikan bahwa berkas yang ia pasang benar-benar berasal dari penerbit yang
dipercaya dan tidak berubah sejak ditandatangani.

\textit{Digital signature} menjawab kebutuhan tersebut. Penerbit menandatangani sebuah pesan
dengan kunci privatnya, dan siapa pun yang memiliki kunci publik yang sesuai dapat memeriksa
tanda tangan itu. Dengan begitu, penerima dapat memastikan dua hal, yakni \textit{authenticity} (tanda
tangan dibuat oleh pemilik kunci privat) dan \textit{integrity} (pesan tidak diubah sesudah
ditandatangani). Tanda tangan digital tidak menyatakan bahwa perangkat lunak bebas dari malware
atau kerentanan. Yang dijamin hanyalah asal dan keutuhan berkas, sedangkan kepercayaan terhadap
penerbit ditentukan oleh kunci publik mana yang diterima pengguna.

Nadina dirancang sebagai kakas baris perintah yang dijalankan di dalam pipeline CI/CD untuk
menandatangani artefak hasil build dan memverifikasinya sebelum publikasi atau deployment.
Program didistribusikan sebagai image container agar pipeline tidak perlu memasang
dependensinya sendiri. Dua tahap dipisahkan, yakni penandatanganan yang memerlukan
kunci privat, kemudian tahap verifikasi yang hanya membutuhkan artefak, bundle, dan trust store,
sehingga verifikasi dapat dijalankan tanpa menyentuh kunci penerbit. Pengelolaan \textit{trusted public key} 
dilakukan pengguna lewat trust store yang terpisah dari artefak. Konsep alur
penandatanganan ini mengikuti Sigstore Cosign, dengan algoritma dan format yang ditetapkan
spesifikasi tugas.

\subsection{Tujuan dan Ruang Lingkup}
Program ini harus memenuhi beberapa tujuan berikut:
\begin{itemize} 
    \item Program yang dirancang harus dapat melakukan proses penandatanganan dan verifikasi
rilis perangkat lunak (software release) menggunakan algoritma RSA, ECDSA, dan EdDSA.
    \item Program ditujukan untuk digunakan sebagai kakas pada pipeline CI/CD. Program
didistribusikan sebagai image container yang dapat diambil dari registry, kemudian
dijalankan untuk menandatangani artefak hasil build dan memverifikasinya sebelum
publikasi atau deployment.
    \item Program wajib menyediakan pembangkitan pasangan kunci,
pembentukan release manifest, dan pengelolaan kunci publik tepercaya
    \item Pengujian dan analisis dilakukan pada ketiga algoritma untuk membandingkan kebenaran hasil, keamanan
sistem, dan performa program.
    \item Program yang dibuat mengadaptasi alur penandatanganan berkas pada Sigstore Cosign.
\end{itemize}
Berkas rilis yang akan ditandatangani selanjutnya disebut artefak. Informasi mengenai
artefak, seperti nama paket, versi, ukuran berkas, dan nilai hash, disimpan dalam release
manifest. Tanda tangan digital diberikan pada manifest tersebut, kemudian manifest dan
tanda tangannya disimpan dalam satu berkas yang disebut bundle.
